\chapter{Conclusion}
\label{ch:conclusion}
This project set out to apply U-Net-based RFI detection, first by reproducing
the method of Akeret et al.\ and then by designing an improved network, and to
test both on real telescope data rather than only on synthetic data.

The authors' \tfunet{} first scored an \Fone{} of 0.39 on my synthetic data.
Controlled experiments showed that this was caused by two training settings,
class weighting combined with the network's ReLU output layer, and per-image
normalisation, and not by its architecture. Correcting them raised the score to
0.93 without changing the network.

On real LOFAR data labelled by a human expert, the corrected \tfunet{} reached
a maximum \Fone{} of $0.5482 \pm 0.0139$, close to the U-Net result published for
the same data. My network, \hybrid{}, reached $0.6603 \pm 0.0040$ with under
600{,}000 parameters, above the best published result of 0.5979. Independent
checks found no leakage or evaluation error behind this result.

Removing the strip convolutions, channel attention and residual blocks that
\hybrid{} was designed around did not measurably change its score, and adding
a CRF refinement did not improve it. The advantage over \tfunet{} therefore
comes from other differences in the network and its training, which remain to
be identified.

Along the way, checking the public HERA and LOFAR datasets showed that both
contain test images that also appear in the training set, and comparing the
same network on synthetic and real data showed a drop of 0.44 in \Fone{}. Both
point to the same lesson as the rest of this work: results need to be checked
on real data, with repeated runs, before their explanations can be trusted.
